\section{The translation: rainbow trees in cyclic colourings}\label{sec:translation}

\begin{definition}\label{def:colourings}
\leavevmode
\begin{enumerate}
\item The \term{near-distance (ND) colouring} of $K_{2n+1}$ on $\Z_{2n+1}$ gives the edge $xy$ the colour $\ell(x,y)\in\{1,\dots,n\}$.
\item The \term{graceful colouring} of $K_{n+1}$ on $\{0,\dots,n\}$ gives $xy$ the colour $|x-y|$.
\item For an abelian group $G$, the \term{sum colouring} $K_G$ gives $xy$ the colour $x+y$~\cite[Definition~1.2]{mp25b}.
\end{enumerate}
\end{definition}

\begin{lemma}\label{lem:colourings}
\leavevmode
\begin{enumerate}
\item In the ND colouring each colour class is a $2$-regular spanning graph, and every translation $x\mapsto x+g$ preserves all colours.
\item In $K_G$ each colour class is a matching, and $x\mapsto x+g$ maps colour $c$ to colour $c+2g$.
\end{enumerate}
\end{lemma}

\begin{proof}
For (i), $\ell(x+g,y+g)=\ell(x,y)$; for (ii), $(x+g)+(y+g)=c+2g$.
\end{proof}

\begin{proposition}\label{prop:kotzig}
Let $T$ be a tree with $n$ edges. Then $K_{2n+1}$ has a cyclic decomposition into copies of $T$ if and only if the ND-coloured $K_{2n+1}$ contains a rainbow copy of $T$, if and only if $T$ has a $\rho$-labelling.
\end{proposition}

\begin{proof}
A rainbow copy is the same thing as a $\rho$-labelling. A rainbow copy gives a cyclic decomposition by the argument of Theorem~\ref{thm:rosa}, taking its $2n+1$ translates and using Lemma~\ref{lem:colourings}(i). The converse requires checking that all copies form a single orbit. Suppose a decomposition of $K_{2n+1}$ into copies of $T$ is invariant under translation in $\Z_{2n+1}$.
\begin{enumerate}
\item No non-zero translation fixes a copy. Otherwise it induces a fixed-point-free automorphism of odd order of a tree. But a tree automorphism fixes a vertex or an edge, and a map of odd order cannot swap the ends of an edge.
\item Hence each orbit of copies has size $2n+1$. There are exactly $2n+1$ copies, so they form one orbit $T_0+i$.
\item Each length class has $2n+1$ edges, and translation preserves lengths, so $T_0$ contains exactly one edge of each length.
\end{enumerate}
So $T_0$ is rainbow in the ND colouring; this is the cyclic formulation of~\cite[\S1]{mps20}.
\end{proof}

\begin{proposition}\label{prop:graceful}
For a tree $T$ with $n$ edges the following are equivalent:
\begin{enumerate}
\item $T$ has a graceful labelling;
\item the graceful colouring of $K_{n+1}$ has a rainbow spanning copy of $T$;
\item the ND-coloured $K_{2n+1}$ has a rainbow copy of $T$ whose vertex set is $\{0,1,\dots,n\}$.
\end{enumerate}
\end{proposition}

\begin{proof}
A bijection $V(T)\to\{0,\dots,n\}$ is a spanning copy, and ``distinct $|f(x)-f(y)|$'' means ``rainbow'', so (i) and (ii) are equivalent. For $x,y\in\{0,\dots,n\}$ we have $|x-y|\leq n$, so $\ell(x,y)=|x-y|$, and (ii) and (iii) are equivalent.
\end{proof}

Fixing the vertex set ``adds an additional restriction which makes it harder to find such a rainbow copy''~\cite[\S1]{mps20}.

\begin{proposition}[M\"uyesser--Pokrovskiy~\cite{mp25b}, \S1]\label{prop:harm}
Let $T$ be a tree with $n$ edges and let $T'=T-v$ for a leaf $v$. If $K_{\Z_n}$ contains a rainbow copy of $T'$ on all $n$ vertices, then $T$ is harmonious.
\end{proposition}

\begin{proof}
The copy $\varphi\colon V(T')\to\Z_n$ is a bijection with $n-1$ distinct edge sums, so exactly one colour $c$ is unused. Let $w$ be the neighbour of $v$ and set $\varphi(v)=c-\varphi(w)$. Now the $n$ edge sums are all of $\Z_n$, and exactly one label is repeated.
\end{proof}

The converse can fail, since the repeated label need not sit on a leaf, so the rainbow statement is stronger than harmoniousness.

\begin{proposition}[Beals et al.~\cite{bghj91}, \S2]\label{prop:hgroup}
Colour the arc $x\to y$, $x\neq y$, of the complete digraph on $G$ by $xy$. Then $G$ is harmonious if and only if this digraph has a rainbow directed Hamilton cycle. For abelian $G$, this holds if and only if $K_G$ has a rainbow Hamilton cycle.
\end{proposition}

\begin{proof}
A harmonious list $g_1,\dots,g_n$ is the Hamilton cycle $g_1\to g_2\to\dots\to g_n\to g_1$, whose $n$ arc colours $g_ig_{i+1}$ are all of $G$ and hence distinct; conversely, a rainbow Hamilton cycle has $n$ distinct colours from a set of size $n$. For abelian $G$ we have $xy=yx$, so arcs can be replaced by edges.
\end{proof}

Compare the Cayley digraphs of Graham's problem: there the arc $x\to xs$ is coloured $s$, a \emph{difference}; here $x\to y$ is coloured $xy$, a \emph{product}.

\begin{proposition}[Ollis--Pasotti--Pellegrini--Schmitt~\cite{opps20}, \S1]\label{prop:bhr}
For a multiset $L$ of $v-1$ elements of $\{1,\dots,\lfloor v/2\rfloor\}$ the following are equivalent:
\begin{enumerate}
\item $K_v$ has a Hamilton path with edge lengths $L$;
\item $L=\{a_1,\dots,a_{v-1}\}$ can be ordered and given signs $\varepsilon_i=\pm1$ so that the partial sums of $\varepsilon_1a_1,\dots,\varepsilon_{v-1}a_{v-1}$ are exactly $\Z_v\setminus\{0\}$;
\item the circulant multigraph $\operatorname{Cay}[\Z_v:L\cup-L]$ has a cyclic decomposition into Hamilton paths~\cite{pp14}.
\end{enumerate}
\end{proposition}

\begin{proof}[Proof of the equivalence of (i) and (ii)]
Translate the path so that it starts at $0$. Its steps $x_{i+1}-x_i=\pm a_i$ have lengths $a_i$, and its vertices $0,x_1,\dots,x_{v-1}$ are the partial sums; ``Hamilton'' means that they are all of $\Z_v$.
\end{proof}

For $p=7$, the steps $1,1,3,-2,3,-2$ have partial sums $1,2,5,3,6,4$, which is the path of Figure~\ref{fig:buratti}.

\topic{The dictionary}

\begin{center}
\renewcommand{\arraystretch}{1.25}
\small
\begin{tabular}{p{0.44\textwidth}p{0.48\textwidth}}
\textbf{labelling / group} & \textbf{rainbow subgraph}\\\hline
$\rho$-labelling of $T$ & rainbow $T$ in the ND-coloured $K_{2n+1}$\\
graceful labelling of $T$ & rainbow spanning $T$ in the graceful colouring of $K_{n+1}$\\
harmonious labelling of $T$ & $\Leftarrow$ rainbow $T-v$ in $K_{\Z_n}$\\
harmonious group $G$ & rainbow directed Hamilton cycle, arc $x\to y$ coloured $xy$\\
Buratti / Horak--Rosa & Hamilton path with prescribed colour multiset\\
cyclic decomposition of $K_{2n+1}$ & one rainbow copy, translated\\
Maamoun--Meyniel & no rainbow Hamilton path in $K_{\mathbb F_2^p}$
\end{tabular}
\end{center}
